\section{Cumulative distribution functions}

The cumulative distribution function is the most universal way to describe a
real-valued random variable. It works for discrete random variables, continuous
random variables, and more complicated random variables. In the continuous case,
it is especially important because probabilities are naturally built from
intervals.

\begin{definitionbox}
For a real-valued random variable \(X\), the \textbf{cumulative distribution
function}, or \textbf{CDF}, is the function \(F_X:\mathbb{R}\to[0,1]\) defined by
\[
F_X(x)=P(X\leq x).
\]
\end{definitionbox}

The value \(F_X(x)\) is the probability accumulated up to the point \(x\). It
answers the question: how much probability lies to the left of \(x\)?

\subsection*{CDF from a density}

If \(X\) has density \(f_X\), then
\[
F_X(x)=P(X\leq x)=\int_{-\infty}^{x}f_X(t)\,dt.
\]
The variable \(t\) is a dummy integration variable. It is used so that \(x\) can
remain the endpoint of accumulation.

This formula says that the CDF is an accumulated area function. As \(x\) moves to
the right, more area under the density is included.

\subsection*{Recovering interval probabilities from the CDF}

For \(a<b\), the event \(X\leq b\) can be split into two parts:
\[
\{X\leq b\}=\{X\leq a\}\cup\{a<X\leq b\}.
\]
These two events are disjoint. Therefore
\[
P(X\leq b)=P(X\leq a)+P(a<X\leq b).
\]
Rearranging gives
\[
P(a<X\leq b)=F_X(b)-F_X(a).
\]
For continuous random variables with no point masses, the endpoint convention is
irrelevant, so we also write
\[
P(a\leq X\leq b)=F_X(b)-F_X(a).
\]

\begin{theorembox}
For a continuous random variable with CDF \(F_X\), interval probabilities are
computed by differences:
\[
P(a\leq X\leq b)=F_X(b)-F_X(a).
\]
\end{theorembox}

\subsection*{Example: uniform distribution on \([0,1]\)}

Let \(X\) be uniformly distributed on \([0,1]\). Its density is
\[
f_X(x)=
\begin{cases}
1, & 0\leq x\leq 1,\\
0, & \text{otherwise}.
\end{cases}
\]
For \(x<0\), no probability lies to the left of \(x\), so
\[
F_X(x)=0.
\]
For \(0\leq x\leq 1\),
\[
F_X(x)=\int_0^x1\,dt=x.
\]
For \(x>1\), all probability lies to the left of \(x\), so
\[
F_X(x)=1.
\]
Thus
\[
F_X(x)=
\begin{cases}
0, & x<0,\\
x, & 0\leq x\leq 1,\\
1, & x>1.
\end{cases}
\]
For example,
\[
P(0.2\leq X\leq 0.7)=F_X(0.7)-F_X(0.2)=0.7-0.2=0.5.
\]

\subsection*{Example: triangular density}

Let
\[
f_X(x)=
\begin{cases}
2x, & 0\leq x\leq 1,\\
0, & \text{otherwise}.
\end{cases}
\]
For \(0\leq x\leq1\),
\[
F_X(x)=\int_0^x2t\,dt=x^2.
\]
Therefore
\[
F_X(x)=
\begin{cases}
0, & x<0,\\
x^2, & 0\leq x\leq 1,\\
1, & x>1.
\end{cases}
\]
Then
\[
P(0.5\leq X\leq 1)=F_X(1)-F_X(0.5)=1-0.25=0.75.
\]
This agrees with the direct integral calculation.

\subsection*{Basic properties of every CDF}

Every CDF has the following properties:
\begin{itemize}
    \item \(0\leq F_X(x)\leq1\) for every \(x\);
    \item \(F_X\) is non-decreasing;
    \item \(\lim_{x\to-\infty}F_X(x)=0\);
    \item \(\lim_{x\to\infty}F_X(x)=1\).
\end{itemize}

For continuous random variables with densities, the CDF is usually continuous.
If the density is continuous at \(x\), then the derivative of the CDF at \(x\) is
the density:
\[
F_X'(x)=f_X(x).
\]
This is an application of the fundamental theorem of calculus.

\begin{remarkbox}
The density is the rate at which the CDF accumulates probability. A high density
means that the CDF rises quickly. A low density means that the CDF rises slowly.
\end{remarkbox}

\subsection*{CDF versus density}

The CDF and density describe the same continuous distribution, but they answer
different questions.

The density answers a local question: where is probability concentrated per unit
length? The CDF answers an accumulated question: how much probability has been
collected up to \(x\)?

For computations, the CDF is often convenient because interval probabilities are
simple differences:
\[
F_X(b)-F_X(a).
\]
For interpretation, the density is often convenient because its shape shows where
probability is concentrated.

\begin{summarybox}
For a continuous random variable with density \(f_X\),
\[
F_X(x)=\int_{-\infty}^{x}f_X(t)\,dt,
\qquad
P(a\leq X\leq b)=F_X(b)-F_X(a).
\]
The density gives local concentration; the CDF gives accumulated probability.
\end{summarybox}
